% from ega4-2-fr.tex
\oldpage[IV-2]{151}
\item[(iii)] $(S_n)$ au point $x$;
\item[(iv)] régulier au point $x$;
\item[(v)] $(R_n)$ au point $x$;
\item[(vi)] normal au point $x$;
\item[(vii)] réduit au point $x$;
\end{enumerate}
si $f$ est plat au point $x$, et si en outre la propriété correspondante
$P(f^{-1}(f(x)),x)$ (notation de (6.7.8)) est vraie.

On dit que $f$ est lisse au point $x$ s'il est localement de présentation finie
dans un voisinage de $x$ dans $X$, et s'il est régulier au point $x$. On dit
respectivement que $f$ est un morphisme: de coprofondeur $\leq n$, de
Cohen--Macaulay, $(S_n)$, régulier, $(R_n)$, normal, réduit, lisse, s'il a la
propriété correspondante en tout point de $X$.
\end{definition}

\begin{proposition}[6.8.2]
\label{IV.2.s6.8.2-fr}
Soient $X$, $Y$ deux préschémas localement noethériens, $f:X\to Y$ un
morphisme, $x$ un point de $X$. Désignons par $M(f,x)$ une des propriétés
(i) à (vii) de la définition (6.8.1), ou la propriété pour $f$ d'être lisse au
point $x$. Soient $Y'$ un préschéma localement noethérien, $g:Y'\to Y$ un
morphisme, $X'=X\times_Y Y'$, $f'=f_{(Y')}:X'\to Y'$. On suppose que $f$
ou $g$ est localement de type fini. Alors, pour tout $x'\in X'$ au-dessus de
$x$, la propriété $M(f,x)$ entraîne $M(f',x')$.
\end{proposition}

\begin{proof}
Posons $y=f(x)$, $y'=f'(x')$; par transitivité des fibres (I, 3.6.4), on a
\[
f'^{-1}(y')=f^{-1}(y)\otimes_{k(y)}k(y').
\]
Comme ou bien $f^{-1}(y)$ est localement de type fini sur $k(y)$, ou bien
$k(y')$ est de type fini sur $k(y)$ (I, 6.4.11), il résulte de (6.7.8) que
les propriétés $P(f^{-1}(y),x)$ et $P(f'^{-1}(y'),x')$ sont équivalentes; en
outre, si $f$ est plat au point $x$, $f'$ est plat au point $x'$ (2.1.4), ce
qui prouve la proposition, compte tenu de (1.4.3, (iii)).
\end{proof}

\begin{proposition}[6.8.3]
\label{IV.2.s6.8.3-fr}
Pour un morphisme $f$ de préschémas localement noethériens, soit $M(f)$
l'une des propriétés suivantes: être de Cohen--Macaulay, $(S_n)$, régulier,
$(R_n)$, normal, réduit.
\begin{enumerate}
\item[(i)] Soient $X$, $Y$, $Z$ trois préschémas localement noethériens,
$f:X\to Y$, $g:Y\to Z$ deux morphismes. Si $M(f)$ et $M(g)$ sont vraies,
$M(g\circ f)$ est vraie.
\item[(ii)] Inversement, si $f$ est surjectif et si $M(f)$ et $M(g\circ f)$
sont vraies, alors $M(g)$ est vraie.
\item[(iii)] Soient $X$, $Y$, $Y'$ trois préschémas localement noethériens,
$f:X\to Y$, $h:Y'\to Y$ deux morphismes; on pose $X'=X\times_Y Y'$,
$f'=f_{(Y')}:X'\to Y'$. On suppose que $f$ ou $h$ est localement de type
fini. Alors, si $M(f)$ est vraie, il en est de même de $M(f')$; la réciproque
est vraie lorsque $h$ est fidèlement plat.
\end{enumerate}
Les conclusions de (i) et (iii) sont encore vraies lorsque $M$ est la propriété
d'être lisse et que (dans (iii)) $h$ est quasi-compact.
\end{proposition}

\begin{proof}
